\documentclass[11pt,a4paper]{letter}
\usepackage{fontspec}
\setmainfont{cmunrm.otf}[BoldFont=cmunbx.otf,ItalicFont=cmunti.otf,BoldItalicFont=cmunbi.otf]
\usepackage[margin=2.5cm]{geometry}
\usepackage[hidelinks]{hyperref}
\setlength{\parskip}{0.6em}
\signature{Rostyslav Zatserkovnyi\\Lviv University of Trade\\and Economics\\\href{mailto:zatserkovnyi.rostyslav@gmail.com}{zatserkovnyi.rostyslav@gmail.com}\\ORCID 0000-0001-6991-2866}
\address{Rostyslav Zatserkovnyi\\Lviv University of Trade and Economics\\10 Tuhan-Baranovskyi Str.\\Lviv 79008, Ukraine}
\date{}
\begin{document}
\begin{letter}{The Editor-in-Chief\\Universal Access in the Information Society\\Springer Nature}
\opening{Dear Editor,}

We submit the research article \emph{``Detecting Text-Reading Hallucinations in Vision-Language Descriptions of Cyrillic Signage, Packaging and Documents for Blind Users''} by Rostyslav Zatserkovnyi, Roksoliana Zatserkovna and Petro Kutsyk for consideration in \emph{Universal Access in the Information Society}.

\textbf{Context.} Reading printed text has always been one of the most frequent reasons blind people take photographs: they want to know what a sign, a product label, a medicine box or a letter says. In the last few years, general-purpose vision-language models have taken over much of this task, and scene-description apps now answer a photograph with a fluent spoken description. That fluency is the problem our work addresses. A classical OCR engine that cannot read a word usually produces visible nonsense. A vision-language model tends to produce a \emph{plausible} word instead: a familiar brand in place of an unusual one, a complete street name where the sign is cut off, or a date or amount that is one digit off. A blind user usually has no independent way to notice this, and earlier studies show that blind users tend to trust machine descriptions. A wrong expiry date, dosage or sum due is therefore an accessibility failure, not a cosmetic one. The risk is plausibly higher for Cyrillic and for Ukrainian in particular, whose specific letters (і, ї, є, ґ) and dense proper names a model with a generic Cyrillic prior can silently ``correct''.

\textbf{What the work does.} We measure how often four vision-language models of one vendor (Anthropic's Claude family) assert text that is not printed, and whether such assertions can be flagged at run time without ground truth. A controlled study on 567 synthetic images varies image quality, framing (cropping, finger and glare occlusion) and word type. A second study on 176 real images (Open Food Facts packaging and official Ukrainian documents) repeats the measurements with model-independent ground truth, and a paired probe edits one letter on real document pages. Every quoted word is scored deterministically against character-level visibility ground truth. The labeller is checked by two automated raters from different vendors (Claude Sonnet and the open Gemma model), and the main comparisons use paired tests with multiple-comparison correction.

\textbf{Main findings.} Only 2.8--5.6\% of quoted words are hallucinated, yet 35--73\% of descriptions contain at least one such word, mostly without any hedge. On real documents, a near-word one letter away from a common word was reported as the common word on 44--56\% of pages. Combining the model's own second reading, an independent OCR pass and agreement between models detects hallucinated words with an AUROC of 0.86--0.97, so an assistive application can tell the user \emph{which} words to doubt, at a cost and latency we report per image.

\textbf{Fit with the journal.} The work concerns equitable access to printed information for blind and visually impaired people, through technology that is now in everyday use, and it gives developers of assistive reading applications concrete and inexpensive safeguards. It also concerns a language and script that are under-represented in accessibility research. We believe these points match the journal's focus on universal access, assistive technologies and the inclusion of users with disabilities. We are explicit about the limits of the study: it covers one model family and contains no user study. Whether flagged words actually help blind users is the next question, and we release the benchmarks, ground truth and code to support it.

The manuscript is original, has not been published previously, and is not under consideration elsewhere. All authors have approved the submission and declare no competing interests. The use of generative AI in preparing the work is disclosed in the manuscript's declarations.

Thank you for considering our manuscript.

\closing{Sincerely,}
\end{letter}
\end{document}
